\pdexercisesfor{\S\ref{sec:ls-many-authors}}{Many authors, one sovereign}

\begin{pdexercise}[kind=Check,rating=1]{ex:ls-many-authors-tree}
Bob holds the charter. He delegates to Alice, to Carol, and to the scope
\texttt{tests/}. Alice and Carol each delegate to the scope \texttt{payments/}.
Using Theorem~\ref{thm:ls-unity-of-command}, decide whether escalation is well
defined for a witness touching \texttt{tests/}, and for one touching
\texttt{payments/}. Give two different repairs for the case that fails.
\end{pdexercise}
\begin{pdsolution}{ex:ls-many-authors-tree}
For \texttt{tests/} the principals with a path to it are $\{\text{Bob}\}$, a
chain, so escalation goes to Bob. For \texttt{payments/} they are
$\{\text{Alice},\text{Carol},\text{Bob}\}$. Alice and Carol are both parents of
\texttt{payments/} in the transitive reduction and neither reaches the other,
so the set is not a chain: the first hop could be either of them.

Repair one: remove one of the two edges into \texttt{payments/}, so a single
person owns it. Repair two: let Carol's authority come through Alice, so Alice
delegates to Carol, giving the path Bob, Alice, Carol, \texttt{payments/}. The edge from Alice to
\texttt{payments/} is then implied and drops out of the reduction, and the set
$\{\text{Carol},\text{Alice},\text{Bob}\}$ is a chain with Carol first. A
declared order between Alice and Carol would also work, but it amounts to the
second repair chosen by rule rather than by assignment.
\end{pdsolution}
\begin{pdexercise}[kind=Trace,rating=2]{ex:ls-many-authors-cycle}
In the worked example, the migration takes thirty minutes, Bob's review takes
$R$ minutes and must follow it, and Alice's demo is $D$ minutes away. Write the
cycle weight as a function of $R$ and $D$. Find the longest feasible review
when $D=80$ and when $D=70$, and the earliest feasible demo when Bob insists
on $R=60$.
\end{pdexercise}
\begin{pdsolution}{ex:ls-many-authors-cycle}
With $z$ the present moment, the four constraints are $s-m\le-30$, $m-e\le-R$,
$z-s\le0$ and $e-z\le D$. Summing them, the left side is $0$, so the cycle
weight is $-30-R+0+D$ and the schedule is feasible exactly when $D-30-R\ge0$,
that is $R\le D-30$, with the migration starting now.

At $D=80$ the longest review is $80-30=50$ minutes, with weight
$-30-50+80=0$ and no slack. At $D=70$ it is $70-30=40$ minutes; a fifty-minute
review there gives $-30-50+70=-10$, infeasible. With $R=60$ the demo must be
at least $30+60=90$ minutes away, with weight $-30-60+90=0$ at that point.
\end{pdsolution}
\begin{pdexercise}[kind=Trace,rating=2]{ex:ls-many-authors-bets}
In Critch and Russell's cake example, Alice gives the cake being red
probability $0.9$ and Bob gives it $0.1$. Each values no cake at $0$, half at
$20$ and all of it at $30$. Compute each person's expected utility, by their
own beliefs, under: an even split; ``all to Alice if red, all to Bob if
green''; and ``all to Alice'' regardless of color. Which policy does each
person prefer? How would a float-plan bond make the second policy explicit
instead of leaving it to an agent's learned weights?
\end{pdexercise}
\begin{pdsolution}{ex:ls-many-authors-bets}
Even split: $20$ for each. The color-contingent policy: Alice expects
$0.9\times30+0.1\times0=27$, and Bob, who believes green with probability
$0.9$, expects $0.1\times0+0.9\times30=27$. ``All to Alice'': Alice $30$,
Bob $0$. Both prefer the contingent policy to the split. Alice prefers ``all
to Alice'' most of all, and Bob prefers the contingent policy to it by $27$ to
$0$. It is a bet on the color, which is why no
fixed weighting of their utilities produces it.

As a float plan, each person stakes their half of the cake against the
observable color before the agent acts. The terms appear in the digest, both
consent in advance, and settlement follows the color. The agent then only
executes a contract both people signed; it does not decide whose belief
deserves more weight.
\end{pdsolution}
\begin{pdexercise}[kind=Open,rating=3]{ex:ls-many-authors-pause}
Rule R9 lets a principal who does not own a shared scope pause its
safety-class actions for a declared time-to-live $\tau$. Specify the rule
completely: who may pause, for how long, whether a pause may be renewed, who
must answer it, and what an unupheld pause costs. Bound the longest time the
shared work can stay paused if all $M$ principals act against it in turn, and
say what your bound assumes about the owner's response time.
\end{pdexercise}
\begin{pdsolution}{ex:ls-many-authors-pause}
One workable specification: any principal holding a grant inside the shared
scope may pause its safety-class actions once per incident for $\tau$; the
pause cannot be renewed by the same principal; it escalates to the scope's
owner, who must uphold it (turning it into the owner's own stop) or lift it.
An unupheld pause carries the dismissal debit of
Theorem~\ref{thm:escalation-threshold}, so a principal pauses only when the
expected benefit clears the debit; the band analysis of that theorem then
applies per principal.

If the owner never answers, the work stays paused at most $M\tau$: each of the
$M$ principals can add one non-renewable pause. If the owner answers within
$\omega$ of each pause, an upheld pause becomes the owner's decision and is no
longer a pause, while a lifted one costs at most $\omega$, so the
pause-attributable delay is at most $M\min(\tau,\omega)$. The bound assumes
pauses are serialized by the daemon and that ``incident'' is defined so a
principal cannot open new incidents at will; the second assumption is the
open part, and it is where a real design must say what counts as the same
incident.
\end{pdsolution}
